\section{Orders: Another Kind of Relation}
\chapterline

In Chapter~15, we studied equivalence relations. An equivalence relation inspects a universe of objects and gathers them into subsets of equivalent elements: two fractions like $2/4$ and $3/6$ represent the same numerical value; two integers sharing a remainder modulo $12$ land on the exact same hour of a clock face. Equivalence partitions a set into disjoint equivalence classes.

An \textbf{order relation} does something different. Rather than grouping items into classes of equivalent elements, an order introduces hierarchy: it specifies which elements precede, contain, or are smaller than other elements.

On the real line $\R$, every pair of numbers is comparable: given $x$ and $y$, either $x < y$, $x = y$, or $x > y$. But many mathematical structures do not arrange themselves along a single linear chain.

Consider getting dressed in the morning. You must put on socks before shoes, and a shirt before a sweater. But which comes first: your left sock or your shirt? Your watch or your belt? Neither. You can put your watch on before or after your shirt without any conflict. The morning routine has a definite hierarchy---you cannot put on shoes before socks---yet the tasks do not form a single straight line. They branch. 

Branching hierarchies appear throughout mathematics: in set theory, $\{1, 2\}$ contains $\{1\}$, but between $\{1, 2\}$ and $\{2, 3\}$, neither is contained in the other. To capture these branching hierarchies, mathematics uses a \textbf{partial order}.

\begin{historybox}
For centuries, mathematicians treated ordering as identical to the real number line, where any two numbers can be compared. When George Boole (\emph{The Laws of Thought}, 1854) and Charles Sanders Peirce (1870s--1880s) expressed logic and set theory algebraically, they observed that subset inclusion $\subseteq$ and implication $\implies$ act like orders, but with a key difference: two subsets or propositions can be \emph{incomparable}.

In the 1930s, Garrett Birkhoff developed modern order theory in \emph{Lattice Theory}. Birkhoff showed that partially ordered sets and lattices connect Boolean algebra, projective geometry, and functional analysis.

Order theory also revealed challenges in social choice. In 1785, the Marquis de Condorcet showed that individual transitive preferences can produce cyclical majority outcomes ($A > B > C > A$). In 1951, Kenneth Arrow proved in his Impossibility Theorem that no ranked voting system can convert individual preferences into a single collective order while satisfying a small set of fair conditions.
\end{historybox}

\subsection{Partial Orders: The Axiomatic Architecture}

\begin{definition}[Partial Order]
A relation $\preceq$ on a set $A$ is a \textbf{partial order} if it satisfies three structural conditions:
\begin{enumerate}[label=(\roman*)]
\item \textbf{Reflexivity:} For all $a \in A$,
\[
a \preceq a.
\]
(Every element is comparable to itself.)
\item \textbf{Antisymmetry:} For all $a, b \in A$,
\[
(a \preceq b \text{ and } b \preceq a) \implies a = b.
\]
(If two elements precede each other, they are identical.)
\item \textbf{Transitivity:} For all $a, b, c \in A$,
\[
(a \preceq b \text{ and } b \preceq c) \implies a \preceq c.
\]
(Precedence flows reliably along chains.)
\end{enumerate}
A set $A$ equipped with a partial order $\preceq$ is called a \textbf{partially ordered set}, or \textbf{poset} $(A, \preceq)$.
\end{definition}

\begin{intuitionbox}
Pay close attention to \emph{antisymmetry}. In an equivalence relation, symmetry means that if $a \sim b$, then $b \sim a$. In an order relation, two distinct elements cannot precede each other: if $a \preceq b$ and $b \preceq a$, then $a$ and $b$ must be equal. No two distinct elements can each be less than or equal to the other.
\end{intuitionbox}

\begin{marginnotebox}
\textbf{Weak versus Strict Orders:} The symbol $\preceq$ represents a \textbf{weak} partial order (analogous to $\le$ or $\subseteq$), where every element is comparable to itself ($a \preceq a$). In contrast, a \textbf{strict} partial order $\prec$ (analogous to $<$ or $\subset$) is defined by $a \prec b \iff (a \preceq b \text{ and } a \neq b)$. It is irreflexive ($a \not\prec a$) and transitive. Mathematicians use ``partial order'' to mean a weak order unless specified otherwise.
\end{marginnotebox}

\begin{recoverybox}
\textbf{Pause and think:} Is the strict inequality relation $<$ on $\R$ a partial order under our definition? Test the three axioms! Reflexivity fails immediately: $5 < 5$ is false. The relation $<$ is a \emph{strict} order. To be a partial order in our standard sense, reflexivity is mandatory, which is why we work with $\le$.
\end{recoverybox}

\subsection{Partial Versus Total Orders}

\begin{definition}[Comparability and Total Order]
Let $(A, \preceq)$ be a poset. Two elements $x, y \in A$ are \textbf{comparable} if $x \preceq y$ or $y \preceq x$. If neither condition holds, $x$ and $y$ are \textbf{incomparable}, written $x \parallel y$.

A partial order $\preceq$ is a \textbf{total order} (or \textbf{linear order}) if every pair of elements is comparable:
\[
(\forall x, y \in A)(x \preceq y \text{ or } y \preceq x).
\]
\end{definition}

On the real numbers $(\R, \le)$, comparability always holds: given $3$ and $\pi$, we know $3 \le \pi$. The real numbers line up in a single unbroken queue with no branching.

Now consider the power set $\pow(\{1, 2, 3\})$ ordered by subset inclusion $\subseteq$. 
Is $\{1\} \subseteq \{2\}$? No.
Is $\{2\} \subseteq \{1\}$? No.
Neither set contains the other. They are incomparable under inclusion: $\{1\} \parallel \{2\}$. The elements do not line up in a single linear chain; they form a partially ordered set.

Similarly, consider the positive integers $\N = \{1, 2, 3, \dots\}$ ordered by divisibility ($a \mid b$). Is $4 \mid 6$? No. Is $6 \mid 4$? No. Thus $4 \parallel 6$. Divisibility allows incomparable elements that do not divide each other. Divisibility is a genuine partial order, not a total one.

\subsection{Visualizing Posets: Hasse Diagrams}

If we drew every related pair as an arrow in a directed graph, the diagram would quickly become cluttered. For even a small set, reflexive self-loops ($x \preceq x$) appear at every vertex, and transitivity adds many redundant shortcut arrows: if $1 \mid 2$ and $2 \mid 4$, transitivity forces an arrow from $1$ to $4$, from $1$ to $8$, from $1$ to $16$, and so on.

A \textbf{Hasse diagram} (named after the German number theorist Helmut Hasse) strips away this visual noise using three simple rules:
\begin{enumerate}[label=\arabic*.]
\item \textbf{Upward Flow:} If $x \prec y$, place $y$ strictly higher on the page than $x$. Gravity carries the orientation, so arrowheads are unnecessary.
\item \textbf{Omit Reflexive Self-Loops:} Reflexivity ($x \preceq x$) is understood implicitly; no loops are drawn.
\item \textbf{Omit Transitive Shortcuts:} Draw an edge between $x$ and $y$ only if $y$ \emph{covers} $x$.
\end{enumerate}

\begin{definition}[Covering Relation]
Let $(A, \preceq)$ be a poset. We say that $y$ \textbf{covers} $x$ (written $x \lessdot y$) if $x \prec y$ and there is no intermediate element $z \in A$ such that $x \prec z \prec y$.
\end{definition}

In plain words, $y$ covers $x$ if $y$ is an immediate next rung above $x$ on the ladder, with no intervening elements between them. A Hasse diagram draws only the covering edges, letting your eyes reconstruct the full partial order by tracing upward paths.

\begin{figure}[htbp]
\centering
\begin{tikzpicture}[scale=0.92, >=Latex]
  % Panel 1: Divisibility Poset D(12)
  \begin{scope}[xshift=0cm]
    \node[anchor=south, font=\footnotesize\bfseries] at (2.4, 4.3) {1. Divisibility: $D(12)$ under $a \mid b$};
    \draw[rounded corners=4pt, fill=gray!3, draw=gray!25] (-0.3, -1.15) rectangle (5.1, 4.3);
    
    % Level indicators on left
    \node[font=\tiny, text=gray!55, anchor=west] at (0.0, 0.2) {Level 0};
    \node[font=\tiny, text=gray!55, anchor=west] at (0.0, 1.4) {Level 1};
    \node[font=\tiny, text=gray!55, anchor=west] at (0.0, 2.6) {Level 2};
    \node[font=\tiny, text=gray!55, anchor=west] at (0.0, 3.8) {Level 3};
    
    % Nodes
    \node[circle, draw=sourceblue, fill=sourceblue!20, inner sep=2.5pt, font=\footnotesize\bfseries] (n1) at (2.4, 0.2) {$1$};
    \node[circle, draw=targetgreen, fill=targetgreen!20, inner sep=2.5pt, font=\footnotesize\bfseries] (n2) at (1.4, 1.4) {$2$};
    \node[circle, draw=targetgreen, fill=targetgreen!20, inner sep=2.5pt, font=\footnotesize\bfseries] (n3) at (3.4, 1.4) {$3$};
    \node[circle, draw=quotientorange, fill=quotientorange!20, inner sep=2.5pt, font=\footnotesize\bfseries] (n4) at (1.4, 2.6) {$4$};
    \node[circle, draw=quotientorange, fill=quotientorange!20, inner sep=2.5pt, font=\footnotesize\bfseries] (n6) at (3.4, 2.6) {$6$};
    \node[circle, draw=structurepurple, fill=structurepurple!20, inner sep=2.5pt, font=\footnotesize\bfseries] (n12) at (2.4, 3.8) {$12$};
    
    % Covering edges (upward flow)
    \draw[thick, gray!70] (n1) -- (n2);
    \draw[thick, gray!70] (n1) -- (n3);
    \draw[thick, gray!70] (n2) -- (n4);
    \draw[thick, gray!70] (n2) -- (n6);
    \draw[thick, gray!70] (n3) -- (n6);
    \draw[thick, gray!70] (n4) -- (n12);
    \draw[thick, gray!70] (n6) -- (n12);
    
    \node[below, font=\scriptsize, align=center, text=gray!80!black] at (2.4, -0.2) {$2 \nmid 3$ and $3 \nmid 2$ ($2 \parallel 3$);\\no edge is drawn between them};
  \end{scope}

  % Panel 2: Power Set Lattice P({a, b})
  \begin{scope}[xshift=6.0cm]
    \node[anchor=south, font=\footnotesize\bfseries] at (2.4, 4.3) {2. Subset Lattice: $\pow(\{a, b\})$ under $\subseteq$};
    \draw[rounded corners=4pt, fill=gray!3, draw=gray!25] (-0.3, -1.15) rectangle (5.1, 4.3);
    
    % Level indicators on left
    \node[font=\tiny, text=gray!55, anchor=west] at (0.0, 0.2) {Level 0};
    \node[font=\tiny, text=gray!55, anchor=west] at (0.0, 2.0) {Level 1};
    \node[font=\tiny, text=gray!55, anchor=west] at (0.0, 3.8) {Level 2};
    
    % Nodes
    \node[draw=sourceblue, fill=sourceblue!15, rounded corners=2pt, font=\scriptsize\bfseries, inner sep=3pt] (p0) at (2.4, 0.2) {$\emptyset$};
    \node[draw=targetgreen, fill=targetgreen!15, rounded corners=2pt, font=\scriptsize\bfseries, inner sep=3pt] (pa) at (1.35, 2.0) {$\{a\}$};
    \node[draw=targetgreen, fill=targetgreen!15, rounded corners=2pt, font=\scriptsize\bfseries, inner sep=3pt] (pb) at (3.45, 2.0) {$\{b\}$};
    \node[draw=structurepurple, fill=structurepurple!15, rounded corners=2pt, font=\scriptsize\bfseries, inner sep=3pt] (pab) at (2.4, 3.8) {$\{a, b\}$};
    
    % Incomparable indicator
    \draw[<->, thick, trapred!80, dashed] (pa) -- node[above=2pt, font=\tiny, text=trapred!90!black] {$\{a\} \parallel \{b\}$} (pb);
    
    % Edges
    \draw[thick, gray!70] (p0) -- (pa);
    \draw[thick, gray!70] (p0) -- (pb);
    \draw[thick, gray!70] (pa) -- (pab);
    \draw[thick, gray!70] (pb) -- (pab);
    
    \node[below, font=\scriptsize, align=center, text=gray!80!black] at (2.4, -0.2) {Branching reflects incomparability;\\edges show immediate containment};
  \end{scope}
\end{tikzpicture}
\caption{Hasse diagrams reduce visual clutter by omitting reflexive self-loops ($x \preceq x$) and redundant transitive shortcuts. \textbf{Left}: The divisibility poset $D(12)$ under $a \mid b$. Because $2 \nmid 3$ and $3 \nmid 2$, the numbers $2$ and $3$ stand side by side without an edge ($2 \parallel 3$). \textbf{Right}: The power set lattice $\pow(\{a, b\})$ under $\subseteq$. Incomparable singletons $\{a\}$ and $\{b\}$ branch outward and reunite at their join $\{a, b\}$.}
\label{16:fig-hasse-diagrams}
\end{figure}

\subsection{Worked Proof: Divisibility as a Poset}

Let $\N = \{1, 2, 3, \dots\}$ be the positive integers. We define the relation $a \preceq b \iff a \mid b$ ($a$ divides $b$).

\begin{theorem}
Divisibility is a partial order on $\N$.
\end{theorem}

\begin{proof}
We verify the three required conditions:

\textbf{1. Reflexivity:} For any $a \in \N$, $a = a \cdot 1$. Since $1 \in \N$, $a \mid a$, so $a \preceq a$.

\textbf{2. Antisymmetry:} Suppose $a \mid b$ and $b \mid a$.
Then there exist positive integers $k, \ell \in \N$ such that $b = k a$ and $a = \ell b$.
Substituting $b$ into the second equation:
\[
a = \ell (k a) = (\ell k) a.
\]
Since $a > 0$, we divide by $a$ to obtain $\ell k = 1$.
Because $\ell$ and $k$ are positive integers, the only solution to $\ell k = 1$ is $\ell = k = 1$.
Substituting $k = 1$ into $b = k a$ yields $b = a$. Hence antisymmetry holds.

\textbf{3. Transitivity:} Suppose $a \mid b$ and $b \mid c$.
Then $b = j a$ and $c = m b$ for some $j, m \in \N$.
Substituting $b$:
\[
c = m (j a) = (m j) a.
\]
Since $m j$ is a positive integer, $a \mid c$. Hence transitivity holds.

Since all three conditions hold, $(\N, \mid)$ is a poset.
\end{proof}

\begin{trapbox}
Notice why we restricted divisibility to positive integers $\N$ rather than all integers $\Z$! 
On $\Z$, $2 \mid (-2)$ (since $-2 = (-1) \cdot 2$) and $(-2) \mid 2$ (since $2 = (-1) \cdot (-2)$). Yet $2 \neq -2$. Antisymmetry fails on $\Z$ because the equation $\ell k = 1$ has two solutions in $\Z$: $\ell = k = 1$ and $\ell = k = -1$. The presence of invertible units ($\pm 1$) means that mutual divisibility in $\Z$ only guarantees $a = \pm b$ (they are associates), not equality. Mathematical definitions require careful domain awareness.
\end{trapbox}

\subsection{Extremal Elements: Maximal, Minimal, Greatest, and Least}

In everyday language, ``maximal'' and ``maximum'' are often used interchangeably. In order theory, they have distinct mathematical meanings.

\begin{definition}[Extremal Elements]
Let $(A, \preceq)$ be a poset, and let $S \subseteq A$ be a subset.
\begin{enumerate}[label=(\roman*)]
\item An element $M \in S$ is the \textbf{greatest element} (or \textbf{maximum}) of $S$ if for every $x \in S$,
\[
x \preceq M.
\]
(Plain meaning: $M$ is greater than or equal to every element in the set.)
\item An element $m \in S$ is a \textbf{maximal element} of $S$ if for all $x \in S$,
\[
m \preceq x \implies x = m.
\]
(Plain meaning: no element in the set is strictly greater than $m$.)
\item An element $L \in S$ is the \textbf{least element} (or \textbf{minimum}) of $S$ if for every $x \in S$,
\[
L \preceq x.
\]
(Plain meaning: $L$ is less than or equal to every element in the set.)
\item An element $\ell \in S$ is a \textbf{minimal element} of $S$ if for all $x \in S$,
\[
x \preceq \ell \implies x = \ell.
\]
(Plain meaning: no element in the set is strictly less than $\ell$.)
\end{enumerate}
\end{definition}

Notice the logical distinction between these conditions:
\begin{itemize}
\item A \textbf{greatest element} must be comparable to, and greater than or equal to, every element in the set ($x \preceq M$ for all $x$).
\item A \textbf{maximal element} only requires that no element in the set is strictly greater than it ($m \preceq x \implies x = m$).
\end{itemize}

Every greatest element is automatically maximal. In a poset that is not totally ordered, a maximal element need not be comparable to other elements; it only needs to have no element strictly above it.

\begin{theorem}[Uniqueness of Greatest and Least Elements]
If a subset $S$ of a poset has a greatest element, that element is unique. Similarly, if $S$ has a least element, it is unique.
\end{theorem}

\begin{proof}
Suppose $M_1, M_2 \in S$ are both greatest elements of $S$. Since $M_1 \in S$ and $M_2$ is a greatest element, $M_1 \preceq M_2$. Reversing roles, since $M_2 \in S$ and $M_1$ is a greatest element, $M_2 \preceq M_1$. By the antisymmetry of $\preceq$, $M_1 = M_2$. The argument for the least element is identical with inequalities reversed.
\end{proof}

While a poset can have at most \emph{one} greatest element, it can easily have \emph{many} maximal elements!

\begin{example}[Maximal Without Maximum]
Let $S = \pow(\{1, 2, 3\}) \setminus \{\emptyset, \{1, 2, 3\}\}$ be the collection of all non-empty proper subsets of $\{1, 2, 3\}$, ordered by $\subseteq$:
\[
S = \big\{ \{1\}, \, \{2\}, \, \{3\}, \, \{1, 2\}, \, \{1, 3\}, \, \{2, 3\} \big\}.
\]
\begin{itemize}
\item The maximal elements of $S$ are $\{1, 2\}$, $\{1, 3\}$, and $\{2, 3\}$. None of these subsets is contained in any other element of $S$. Yet none of them is a greatest element, because $\{1, 2\} \nsupseteq \{2, 3\}$. The set $S$ has three maximal elements and \emph{no} greatest element.
\item Symmetrically, the minimal elements of $S$ are $\{1\}$, $\{2\}$, and $\{3\}$. None contains any other element of $S$, but none is a least element. The set $S$ has three minimal elements and \emph{no} least element.
\end{itemize}
\end{example}

\begin{figure}[htbp]
\centering
\begin{tikzpicture}[scale=0.92, >=Latex]
  % Panel 1: Unique Greatest Element
  \begin{scope}[xshift=0cm]
    \node[anchor=south, font=\footnotesize\bfseries] at (2.4, 4.3) {1. Unique Greatest Element};
    \draw[rounded corners=4pt, fill=gray!3, draw=gray!25] (-0.3, -1.2) rectangle (5.1, 4.3);
    
    % Nodes
    \node[circle, draw=sourceblue, fill=sourceblue!20, inner sep=2.5pt, font=\footnotesize\bfseries] (b0) at (2.4, 0.6) {$0$};
    \node[circle, draw=gray!80, fill=gray!15, inner sep=2.5pt, font=\footnotesize\bfseries] (ba) at (1.4, 2.0) {$a$};
    \node[circle, draw=gray!80, fill=gray!15, inner sep=2.5pt, font=\footnotesize\bfseries] (bb) at (3.4, 2.0) {$b$};
    \node[circle, draw=targetgreen, fill=targetgreen!25, inner sep=2.5pt, font=\footnotesize\bfseries] (bM) at (2.4, 3.4) {$M$};
    
    % Edges
    \draw[thick, gray!70] (b0) -- (ba);
    \draw[thick, gray!70] (b0) -- (bb);
    \draw[thick, gray!70] (ba) -- (bM);
    \draw[thick, gray!70] (bb) -- (bM);
    
    % Annotations
    \node[above=3pt of bM, font=\scriptsize\bfseries, text=targetgreen!80!black] {Greatest and Maximal};
    \node[below=3pt of b0, font=\scriptsize\bfseries, text=sourceblue!80!black] {Least and Minimal};
    
    \node[below, font=\scriptsize, align=center, text=gray!80!black] at (2.4, -0.25) {Single summit $M$:\\$x \preceq M$ for every element $x$};
  \end{scope}

  % Panel 2: Multiple Maximal Elements
  \begin{scope}[xshift=6.0cm]
    \node[anchor=south, font=\footnotesize\bfseries] at (2.4, 4.3) {2. Multiple Maximal Elements};
    \draw[rounded corners=4pt, fill=gray!3, draw=gray!25] (-0.3, -1.2) rectangle (5.1, 4.3);
    
    % Nodes
    \node[circle, draw=sourceblue, fill=sourceblue!20, inner sep=2.5pt, font=\footnotesize\bfseries] (c0) at (2.4, 0.6) {$0$};
    \node[circle, draw=quotientorange, fill=quotientorange!25, inner sep=2.5pt, font=\footnotesize\bfseries] (cm1) at (1.4, 2.8) {$m_1$};
    \node[circle, draw=quotientorange, fill=quotientorange!25, inner sep=2.5pt, font=\footnotesize\bfseries] (cm2) at (3.4, 2.8) {$m_2$};
    
    % Incomparable indicator at top
    \draw[<->, thick, trapred!80, dashed] (cm1) -- node[above=2pt, font=\tiny, text=trapred!90!black] {$m_1 \parallel m_2$} (cm2);
    
    % Edges
    \draw[thick, gray!70] (c0) -- (cm1);
    \draw[thick, gray!70] (c0) -- (cm2);
    
    % Annotations
    \node[above=3pt of cm1, font=\scriptsize\bfseries, text=quotientorange!80!black] {Maximal};
    \node[above=3pt of cm2, font=\scriptsize\bfseries, text=quotientorange!80!black] {Maximal};
    \node[below=3pt of c0, font=\scriptsize\bfseries, text=sourceblue!80!black] {Least and Minimal};
    
    \node[below, font=\scriptsize, align=center, text=gray!80!black] at (2.4, -0.25) {Two peaks, no greatest element:\\nothing sits above $m_1$, yet $m_1 \parallel m_2$};
  \end{scope}
\end{tikzpicture}
\caption{The vital distinction between a \emph{greatest element} and a \emph{maximal element}. \textbf{Left}: A poset with a single summit $M$. Because $M$ sits above every element ($x \preceq M$ for all $x$), it is the unique greatest element (maximum) and also maximal. \textbf{Right}: A branching poset with two separate peaks $m_1$ and $m_2$. Both $m_1$ and $m_2$ are maximal (no element sits strictly above them), but neither is a greatest element because neither dominates the other ($m_1 \parallel m_2$).}
\label{16:fig-maximal-vs-greatest}
\end{figure}

\begin{recoverybox}
\textbf{Pause and think:} Look at the positive integers with their usual arithmetic order $(\N, \le)$. What is the least element? The number $1$. Is there a maximal element? No: given any $n \in \N$, the number $n+1$ is strictly larger, so no element can be maximal.
\end{recoverybox}

\subsection{Bounds, Suprema, and Lattices}

What happens when a subset of a poset does not contain its own summit, but sits inside a larger universe? We ask what elements in the universe sit above or below it.

\begin{definition}[Bounds, Supremum, and Infimum]
Let $(A, \preceq)$ be a poset, and let $S \subseteq A$.
\begin{enumerate}[label=(\roman*)]
\item An element $u \in A$ is an \textbf{upper bound} for $S$ if $x \preceq u$ for all $x \in S$.
\item If the set of all upper bounds of $S$ has a least element, that element is called the \textbf{least upper bound} or \textbf{supremum} of $S$, written $\sup S$.
\item An element $\ell \in A$ is a \textbf{lower bound} for $S$ if $\ell \preceq x$ for all $x \in S$.
\item If the set of all lower bounds of $S$ has a greatest element, that element is called the \textbf{greatest lower bound} or \textbf{infimum} of $S$, written $\inf S$.
\end{enumerate}
\end{definition}

An upper bound is an element that sits above or at every element of $S$. The supremum is the least among all upper bounds.

When $S$ consists of just two elements $\{a, b\}$, we write:
\[
a \lor b = \sup\{a, b\} \quad (\text{the \textbf{join}}), \qquad a \land b = \inf\{a, b\} \quad (\text{the \textbf{meet}}).
\]
This notation connects familiar operations across different contexts:
\begin{itemize}
\item \textbf{Sets under inclusion ($\subseteq$):} The join is union, and the meet is intersection:
\[
A \lor B = A \cup B, \qquad A \land B = A \cap B.
\]
The union $A \cup B$ is the smallest set that contains both $A$ and $B$; the intersection $A \cap B$ is the largest set contained in both.
\item \textbf{Positive integers under divisibility ($\mid$):} The join is the least common multiple, and the meet is the greatest common divisor:
\[
a \lor b = \operatorname{lcm}(a, b), \qquad a \land b = \gcd(a, b).
\]
The number $\operatorname{lcm}(a, b)$ is the smallest positive integer divisible by both $a$ and $b$; the number $\gcd(a, b)$ is the largest positive integer that divides both.
\item \textbf{Real numbers under standard order ($\le$):} Consider the open interval $S = (0, 1) \subset \R$. Any real number $u \ge 1$ is an upper bound. The lowest upper bound is $\sup S = 1$. Notice that $1 \notin S$: a supremum does not need to belong to the set it bounds! This single observation is central to real analysis (Chapter~20).
\end{itemize}

A poset in which every pair of elements has a meet and a join is called a \textbf{lattice}. In Chapter~31, we will see that lattices provide the algebraic foundation of digital logic circuits, search engines, and set theory.

\begin{center}
\begin{tikzpicture}[scale=0.88, >=Latex]
  % Panel 1: Abstract Meet and Join
  \begin{scope}[xshift=0cm]
    \draw[rounded corners=4pt, fill=gray!3, draw=gray!25] (-2.7, -2.25) rectangle (2.7, 2.6);
    \node[anchor=south, font=\footnotesize\bfseries] at (0, 2.15) {1. Abstract Lattice Diamond};

    % Top: Join (Supremum)
    \node[circle, draw=quotientorange!90!black, fill=quotientorange!25, thick, inner sep=2pt, font=\scriptsize\bfseries] (sup) at (0, 1.35) {$a \lor b$};
    \node[font=\tiny\bfseries, quotientorange!90!black, right=4pt] at (sup.east) {Join: $\sup$};

    % Middle: a and b
    \node[circle, draw=sourceblue, fill=sourceblue!20, thick, inner sep=2.5pt, font=\scriptsize\bfseries] (na) at (-1.2, 0.45) {$a$};
    \node[circle, draw=targetgreen, fill=targetgreen!20, thick, inner sep=2.5pt, font=\scriptsize\bfseries] (nb) at (1.2, 0.45) {$b$};

    % Incomparable dashed link
    \draw[<->, dashed, gray!60, thin] (na) -- node[above=2pt, font=\tiny, text=gray!60] {$a \parallel b$} (nb);

    % Bottom: Meet (Infimum)
    \node[circle, draw=structurepurple, fill=structurepurple!20, thick, inner sep=2pt, font=\scriptsize\bfseries] (inf) at (0, -0.45) {$a \land b$};
    \node[font=\tiny\bfseries, structurepurple!90!black, right=4pt] at (inf.east) {Meet: $\inf$};

    % Covering edges
    \draw[thick, gray!70] (inf) -- (na);
    \draw[thick, gray!70] (inf) -- (nb);
    \draw[thick, gray!70] (na) -- (sup);
    \draw[thick, gray!70] (nb) -- (sup);

    \node[font=\scriptsize, align=center] at (0, -1.65) {Lowest common ceiling (Join)\\Highest common floor (Meet)};
  \end{scope}

  % Panel 2: Concrete Divisibility Lattice {4, 6}
  \begin{scope}[xshift=6.0cm]
    \draw[rounded corners=4pt, fill=gray!3, draw=gray!25] (-2.7, -2.25) rectangle (2.7, 2.6);
    \node[anchor=south, font=\footnotesize\bfseries] at (0, 2.15) {2. Divisibility: $\{4, 6\}$};

    % Top: LCM = 12
    \node[circle, draw=quotientorange!90!black, fill=quotientorange!25, thick, inner sep=2.5pt, font=\scriptsize\bfseries] (lcm) at (0, 1.35) {$12$};
    \node[font=\tiny\bfseries, quotientorange!90!black, right=4pt] at (lcm.east) {$\operatorname{lcm}$};

    % Middle: 4 and 6
    \node[circle, draw=sourceblue, fill=sourceblue!20, thick, inner sep=2.5pt, font=\scriptsize\bfseries] (d4) at (-1.2, 0.45) {$4$};
    \node[circle, draw=targetgreen, fill=targetgreen!20, thick, inner sep=2.5pt, font=\scriptsize\bfseries] (d6) at (1.2, 0.45) {$6$};

    % Incomparable
    \draw[<->, dashed, gray!60, thin] (d4) -- node[above=2pt, font=\tiny, text=gray!60] {$4 \nmid 6$} (d6);

    % Bottom: GCD = 2
    \node[circle, draw=structurepurple, fill=structurepurple!20, thick, inner sep=2.5pt, font=\scriptsize\bfseries] (gcd) at (0, -0.45) {$2$};
    \node[font=\tiny\bfseries, structurepurple!90!black, right=4pt] at (gcd.east) {$\gcd$};

    % Edges
    \draw[thick, gray!70] (gcd) -- (d4);
    \draw[thick, gray!70] (gcd) -- (d6);
    \draw[thick, gray!70] (d4) -- (lcm);
    \draw[thick, gray!70] (d6) -- (lcm);

    \node[font=\scriptsize, align=center] at (0, -1.65) {LCM is the unique join\\GCD is the unique meet};
  \end{scope}
\end{tikzpicture}\\[0.3em]
{\footnotesize\textbf{Meet, Join, and the Geometry of a Lattice:} \textbf{Left}: For any two elements $a, b$ in a lattice, the join $a \lor b$ is the lowest common ceiling (supremum), and the meet $a \land b$ is the highest common floor (infimum). \textbf{Right}: Under integer divisibility, the join of $4$ and $6$ is their least common multiple $\operatorname{lcm}(4, 6) = 12$, and their meet is their greatest common divisor $\gcd(4, 6) = 2$.}
\end{center}

\subsection{Where Orders Matter and Where We Go Next}

Order theory provides the formal vocabulary for hierarchy across mathematics and computer science:
\begin{itemize}
\item \textbf{Compilers and scheduling:} Task and package dependencies form directed acyclic graphs. Topological sorting extends these partial orders into linear execution pipelines.
\item \textbf{Lattices and digital logic:} Meet and join operations form lattices, which provide the algebraic foundation of Boolean logic and set theory.
\item \textbf{Analysis and bounds:} Concepts like suprema, infima, and completeness in real analysis are defined through posets.
\end{itemize}

An order relation describes how elements line up, but it is a static description. On an infinite ordered set like $\N = \{1, 2, 3, \dots\}$, knowing that $1 < 2 < 3 < \dots$ does not by itself prove a statement for all elements at once.

To prove statements across an infinite sequence of ordered steps, we need an inductive principle. In Chapter~17, we turn to \textbf{Induction: Infinite Claims, Finite Proofs}, which uses the order structure of the natural numbers to establish claims for all $n \in \N$ through finite proofs.

\subsection{Exercises}

\begin{enumerate}
\item \textbf{(Mental Check)} Consider the divisibility poset $(\N, \mid)$ on the positive integers $\N = \{1, 2, 3, \dots\}$.
\begin{enumerate}[label=(\alph*)]
\item Determine whether the following pairs are comparable: $(3, 12)$, $(4, 10)$, and $(7, 7)$.
\item Does $(\N, \mid)$ possess a least element? Does it possess a greatest element?
\item Explain why $(\N, \mid)$ is a partial order but not a total order.
\end{enumerate}

\item \textbf{(Subset Inclusion as a Poset)} Let $S$ be any set. Prove that subset inclusion $\subseteq$ is a partial order on the power set $\pow(S)$ by directly verifying reflexivity, antisymmetry, and transitivity.

\item \textbf{(Hasse Diagram of Divisors)} Draw the Hasse diagram for the divisibility poset on the set of divisors of $30$:
\[
D(30) = \{1, 2, 3, 5, 6, 10, 15, 30\}.
\]
\emph{Hint:} Notice that $30 = 2 \cdot 3 \cdot 5$ is a product of three distinct prime numbers. How does your diagram relate to the three-dimensional cube representing $\pow(\{2, 3, 5\})$?

\item \textbf{(Maximal versus Greatest)} Let $S = \pow(\{1, 2, 3\}) \setminus \{\emptyset, \{1, 2, 3\}\}$ be the collection of all proper, non-empty subsets of $\{1, 2, 3\}$, ordered by subset inclusion $\subseteq$.
\begin{enumerate}[label=(\alph*)]
\item List all maximal elements and all minimal elements of $S$.
\item Does $S$ have a greatest element? Does $S$ have a least element?
\item In your own words, explain how a set can have several maximal elements without having any greatest element.
\end{enumerate}

\item \textbf{(The Product Order)} Let $(A, \preceq_A)$ and $(B, \preceq_B)$ be posets. We define the \emph{product order} $\preceq$ on the Cartesian product $A \times B$ by:
\[
(a_1, b_1) \preceq (a_2, b_2) \iff (a_1 \preceq_A a_2 \text{ and } b_1 \preceq_B b_2).
\]
Prove that $\preceq$ is a partial order on $A \times B$.

\item \textbf{(Lexicographic Order versus Product Order)} Let $(A, \le_A)$ and $(B, \le_B)$ be total orders. The \emph{lexicographic order} (or dictionary order) $\le_{\text{lex}}$ on $A \times B$ is defined by:
\[
(a_1, b_1) \le_{\text{lex}} (a_2, b_2) \iff (a_1 <_A a_2) \text{ or } (a_1 = a_2 \text{ and } b_1 \le_B b_2).
\]
\begin{enumerate}[label=(\alph*)]
\item Prove that if $\le_A$ and $\le_B$ are total orders, then $\le_{\text{lex}}$ is a total order on $A \times B$.
\item Find two points in $\R \times \R$ that are comparable under lexicographic order but incomparable under the product order of Exercise~5.
\end{enumerate}

\item \textbf{(Lattice Operations in Divisibility)} In the divisibility poset $(\N, \mid)$, let $a, b \in \N$.
\begin{enumerate}[label=(\alph*)]
\item Prove that the greatest lower bound $\inf\{a, b\}$ is the greatest common divisor $\gcd(a, b)$.
\item Prove that the least upper bound $\sup\{a, b\}$ is the least common multiple $\operatorname{lcm}(a, b)$.
\end{enumerate}
\end{enumerate}

\newpage
